\chapter{相关研究综述}

\section{时序（微）动作定位分析}
\subsection{通用时序动作定位方法}
时序动作定位（Temporal Action Localization, TAL）是视频理解领域的核心任务之一，
旨在从未经剪辑的长视频中识别出所有动作实例的类别，并精确回归其起始与终止时间。
本章从架构设计、特征提取范式以及序列建模机制三个维度对现有的TAL方法进行系统性的综述与分析。

\textbf{架构设计演进}\quad
根据网络架构设计的不同，现有的主流方法通常被划分为三大类~\cite{liu2024end}：
两阶段\cite{chao2018rethinking,lin2019bmn,zeng2019graph,zhao2021video,zhu2021enriching}、
一阶段\cite{zhang2022actionformer, shi2023tridet, yang2024adapting}
以及基于查询\cite{liu2022end}的方法。
早期的两阶段方法遵循“先生成后分类”的范式。
这类方法首先通过滑动窗口或边界采样生成一系列动作候选片段，
随后利用特定的采样模块提取特征进行分类与精修~\cite{chao2018rethinking, lin2019bmn,zhu2021enriching}。
典型的代表如BMN~\cite{lin2019bmn}通过设计双分支网络分别处理起始与终止边界的概率估计，在边界质量上取得了显著提升。
尽管两阶段方法在边界质量上具有优势，但其串行的处理流程限制了推理效率。
相比之下，一阶段方法直接基于多尺度特征金字塔进行端到端的动作定位，极大地简化了检测流程。
以ActionFormer\cite{zhang2022actionformer}为代表的先进模型，
利用Transformer\cite{vaswani2017attention}强大的上下文建模能力，
在单一网络中并行完成动作分类与边界回归，成为了当前性能与效率平衡的标杆。
此外，受目标检测DETR\cite{zhu2020deformable}架构启发，
基于查询的方法（如TadTR~\cite{liu2022end}）通过学习一组稀疏的查询向量与视频特征交互，
实现了无需后处理的直接预测。

\textbf{特征提取范式：从特征基到端到端}\quad
除了架构差异，TAD方法还可根据输入模态的处理方式分为“基于特征”和“端到端”两类。
目前的主流方法（如 ActionFormer）大多属于前者，
即依赖于预训练的视频编码器（如 I3D\cite{carreira2017quo}, VideoMAE\cite{tong2022videomae,wang2023videomae}）离线提取 RGB 或光流特征。
这种解耦策略虽然降低了显存开销，但忽略了特征提取与定位任务的协同优化。
近年来，端到端方法逐渐受到关注。
虽然早期工作如AFSD~\cite{lin2021learning}受限于计算资源需对输入进行大幅下采样，
但后续工作如的DaoTAD~\cite{wang2021rgb}和E2E-TAD~\cite{liu2022empirical}证明了仅凭RGB模态结合强数据增强即可实现优异性能。
为进一步解决显存瓶颈，SGP~\cite{cheng2022stochastic} 和 ETAD~\cite{liu2023etad}提出了部分反向传播策略，
而Re2TAL~\cite{zhao2023re2tal} 则设计了可逆网络架构以释放中间激活占用的显存。
AdaTAD~\cite{liu2024end}进一步将模型的尺寸扩大到10亿参数级别，展示了端到端方法在大规模模型上的潜力。
尽管端到端方法发展迅速，但在考虑模型显存和计算量的情况下，基于特征的方法依旧占据着重要地位。

\textbf{序列建模机制的推演}\quad
时序建模机制经历了从局部感知到全局关联，再到动态选择的选择性演进过程。
早期方法主要依赖一维卷积神经网络（1D-CNN）提取时序特征~\cite{shou2016temporal,lin2017single, xu2020g}，
虽具备局部归纳偏置，但为了扩大感受野而进行的下采样操作往往导致微小动作的细微特征丢失；
随后引入的图卷积网络（GCN）\cite{zeng2019graph,zhao2021video} 尝试将视频片段建模为图节点以聚合非局部语义关系。
为了获取全局上下文，基于Transformer的方法利用自注意力机制~\cite{zhang2022actionformer, yang2024adapting}实现了对长距离依赖的建模，
确立了当前的主流范式，但其二次的计算复杂度限制了高分辨率序列的解析，且静态注意力权重容易被动态的背景“劫持”，难以聚焦微弱的瞬态运动信号。
以 Mamba 为代表的状态空间模型（SSMs）\cite{gu2023mamba, li2024videomamba}
凭借线性复杂度与输入依赖的选择性扫描机制（Selective Scan）建模得到更精细的特征。
不同于传统的静态加权，这种机制赋予了模型内容感知的动态滤波能力，使其能够自适应地过滤冗余的静态背景并敏锐响应微弱的动态变化，
为在强背景干扰下实现精准动作定位提供了全新的技术路径。

\subsection{微动作时序定位研究}



\section{在线动作检测}


\section{医疗场景下的视频动作理解}
\subsection{吞咽造影视频分析}
吞咽造影（Videofluoroscopic Swallowing Study, VFSS）是临床评估吞咽功能的“金标准”，
其核心在于通过分析吞咽过程中的时序逻辑与运动特征来辅助诊断。
早期的研究主要聚焦于半自动分析技术，通过人工干预与计算机算法的结合来提取空间运动学参数。
例如，Kellen等\cite{kellen2010computer}与Hossain等\cite{hossain2014semi}先后探索了基于模板匹配与边缘特征追踪的方法，
利用颈椎等解剖学标志构建相对坐标系，
从而实现舌骨位移、轨迹及活动范围（ROM）的量化。
尽管此类方法在一定程度上提升了客观性并降低了分析耗时，但其高度依赖人工初始化，
且对低对比度或存在解剖重叠的造影图像表现出较强的敏感性，难以实现大规模临床数据的自动化处理。

随着深度学习技术的兴起，自动化时间学参数分析成为了该领域的研究热点方向之一，
其目标是自动识别吞咽各阶段（如口腔期、咽期、食管期）并提取关键时间参数（如咽传输时间、上食管括约肌开合时间等）。
早期的自动化尝试多基于卷积神经网络（CNN）开展。
研究者利用2D或3D CNN对视频帧进行特征提取，通过二分类任务区分咽期与非咽期\cite{lee2019automatic,bandini2021effect,jeong2023application}。
例如，Lee等\cite{lee2019automatic}利用预训练的3D CNN对咽期进行检测，验证了深度学习在VFSS视频分类中的可行性；
Jeong等\cite{jeong2023application}则进一步提出了“从分类到定位”的框架，实现了对剪辑视频中七个不同吞咽阶段的同时识别。
然而，临床实际获取的VFSS数据通常为未剪辑的长视频，且包含大量非典型或微小的动作变化，
这使得简单的分类模型难以满足高精度的临床需求。
针对这一挑战，Ruan等\cite{ruan2023temporal}针对未剪辑视频，
采用多检测器协同的“由粗到细”定位策略，实现了对微吞咽动作的识别与分类。
尽管目前的基于视觉流（RGB）或光流的方法取得了一定进展，但在VFSS场景下仍面临组织对比度低、背景噪声干扰严重等问题。
\subsection{手术工作流识别}

\section{本章小结}